\subsection{A Complete $3\times3$ Example}

Consider
\[
A=\begin{pmatrix}
1&1&0\\0&1&1\\1&0&1
\end{pmatrix}.
\]
Begin with
\[
\left[\begin{array}{ccc|ccc}
1&1&0&1&0&0\\
0&1&1&0&1&0\\
1&0&1&0&0&1
\end{array}\right].
\]

\begin{derivationbox}[Forward and backward elimination]
Apply successively
\[
R_3\leftarrow R_3-R_1,
\qquad
R_3\leftarrow R_3+R_2,
\qquad
R_3\leftarrow \frac12R_3,
\]
then
\[
R_2\leftarrow R_2-R_3,
\qquad
R_1\leftarrow R_1-R_2.
\]
The augmented matrix becomes
\[
\left[\begin{array}{ccc|ccc}
1&0&0&\frac12&-\frac12&\frac12\\
0&1&0&\frac12&\frac12&-\frac12\\
0&0&1&-\frac12&\frac12&\frac12
\end{array}\right].
\]
\end{derivationbox}

Thus
\[
\boxed{A^{-1}=\begin{pmatrix}
\frac12&-\frac12&\frac12\\
\frac12&\frac12&-\frac12\\
-\frac12&\frac12&\frac12
\end{pmatrix}}.
\]

The calculation is longer than in dimension two, but the logic is unchanged: create pivots, eliminate on both sides, and record every row operation on the identity block.
